% =============================================================================
%  empirical_section_updated.tex
%  Numbers synced to output/paper_replication/. Wording fixes applied:
%   - orthogonality footnote: "impact shifts, peak/profile unchanged" (was
%     "virtually unchanged"), + "single interaction term" in the body sentence;
%   - PD-aR terminology unified (was PD-VaR) to match Corollary cor:pd_ar;
%   - rho/Basel sentence added; tab:regZ stars caveat added;
%   - Appendix~\ref (was Section~\ref) for app:alternative_information_sets;
%   - fig:pf caption: shocked September 11 path (was "baseline PD path");
%   - PF "at most about 2%" (was "below 2%").
% =============================================================================
\section{Geopolitical stress testing of U.S. credit risk}
\label{sec:empirical_analysis}
This section applies the framework to U.S. credit risk.\footnote{Code and data:
\url{https://github.com/YoannPull/VAR-MERTON}.} The empirical exercise traces
the transmission of geopolitical-risk innovations from the macro-financial
system to the latent systematic credit factor and to the distribution of
portfolio default probabilities (PDs). After presenting the baseline data and
specification, we use the model for four exercises: a standard impulse-response
analysis, a historical-episode analysis that evaluates large GPR innovations at
their own initial conditions, a comparison with the perfect-foresight
convention used in supervisory stress tests, and a short-sample exercise for
the credit-risk block. We close with robustness checks.
\subsection{Data and empirical specification}
\label{sec:data_empirical_specification}
The application uses quarterly U.S. data from 1986:Q1 to 2024:Q4. The shock
variable is the Geopolitical Risk Index (GPR) of
\citet{caldara2022geopolitical}, aggregated to calendar-quarter frequency.
The index is well suited to the exercise because it provides a long and
event-based measure of geopolitical tensions and realizations, with peaks
corresponding to major geopolitical episodes
(Figure~\ref{fig:GPR_index}).
\begin{figure}[!htbp]
    \centering
    \includegraphics[width=0.80\linewidth]{figures/main/graph_gpr.pdf}
    \caption{\textbf{Geopolitical Risk Index.}
    \textit{Notes: Quarterly averages of the Geopolitical Risk Index,
    1986:Q1--2024:Q4.}}
    \label{fig:GPR_index}
\end{figure}
The baseline credit-risk measure is the Delinquency Rate on All Loans and
Leases at all U.S. commercial banks (FRED: \texttt{DRALACBN}). We use this
long aggregate series as a proxy for realized portfolio-level default risk.
Over the sample, the delinquency rate averages $3.20\%$ and ranges from
$1.19\%$ to $7.50\%$, with the maximum reached during the 2009--2010 credit
downturn (Figure~\ref{fig:US_default}). The Merton--Vasicek inversion of
Section~\ref{sec:macro_credit_architecture} gives a through-the-cycle default
probability $\widehat p=3.20\%$ and an asset correlation
$\widehat\rho=0.051$. This correlation is low relative to the Basel IRB
benchmarks for corporate exposures (between $0.12$ and $0.24$), as expected for
a broadly diversified aggregate delinquency proxy; since the gap between the
quantile and mean responses widens with $\rho$, the tail amplification reported
below should be read as conservative.
\begin{figure}[!htbp]
    \centering
    \includegraphics[width=0.80\linewidth]{figures/main/graph_dralacbn.pdf}
    \caption{\textbf{U.S. aggregate delinquency rate.}
    \textit{Notes: Delinquency rate on all loans and leases, all U.S.
    commercial banks (FRED \texttt{DRALACBN}), quarterly, 1985:Q1--2024:Q4.
    Shaded areas indicate NBER recession periods.}}
    \label{fig:US_default}
\end{figure}
The macro-financial VAR includes the GPR index and five real-side variables:
real private investment per capita, real GDP per capita, private employment
per capita, the real oil price, and year-on-year CPI inflation. This baseline
information set is intentionally parsimonious and focuses on the real-activity
channel through which geopolitical shocks can affect credit risk. Broader
information sets including monetary and uncertainty variables are considered in
Appendix~\ref{app:alternative_information_sets}. Variable definitions and
transformation details are reported in Appendix~\ref{app:variables_detail}.
The VAR is estimated under the Normal--Inverse--Wishart prior described in
Section~\ref{sec:posterior_implementation}. The dynamic experiment is a
one-standard-deviation generalized innovation in the GPR equation. For each
posterior draw, we compute the macro-financial generalized impulse responses
and propagate them through the credit-risk block.
The satellite equation links the reconstructed systematic factor $Z_t$ to
current and lagged macro-financial covariates. To avoid conditioning the
credit-risk transmission on a single selected specification, we average across
parsimonious satellite models using Schwarz weights. Candidate regressors are
built from the non-GPR variables in the baseline information set, with current
values and lags up to four quarters. The implementation of the model averaging
scheme is described in Appendix~\ref{app:satellite_bma}.
The closed-form transmission rests on one key assumption: the satellite error is
orthogonal to the macro-financial innovations ($\eta_t \perp u_t$), so that
geopolitical risk reaches credit risk only through the macro-financial block. We
do not take this restriction for granted but nest it in a one-parameter
control-function family---which reinstates a single interaction term in the
conditional variance---and test it directly; the relaxation and its closed-form
correction are derived in Appendix~\ref{app:orthogonality}.\footnote{Relaxing
the restriction with a control-function term---the standardized reduced-form GPR
innovation entered directly in the satellite---yields a negligible direct
channel ($\widehat\lambda=-0.014$, $90\%$ credible interval
$[-0.063,\,0.034]$; HAC $t=-0.58$, $p=0.56$), and a joint test that the
restricted satellite residual is uncorrelated with the entire reduced-form
innovation vector cannot be rejected ($F=0.24$ on $(6,145)$, $p=0.96$). The
satellite coefficients are unchanged; the peak and the propagation profile of
the PD response are unchanged as well, and only the impact response shifts,
through this statistically insignificant channel; see
Tables~\ref{tab:direct_channel_satellite} and~\ref{tab:direct_channel_pd}.}
Table~\ref{tab:regZ} reports the model-averaged coefficients and posterior
inclusion probabilities for the most relevant satellite terms. The satellite
is estimated on $152$ quarterly observations. The systematic factor loads
positively on real investment, whose contemporaneous term has an inclusion
probability of one, and negatively on the real oil price, with weaker
contributions from inflation and GDP. This pattern gives the credit-risk
transmission a simple economic interpretation: geopolitical shocks that weaken
real activity lower the systematic credit factor and therefore increase
portfolio PDs.
\begin{table}[!htbp]
  \centering
  \scriptsize
  \tablestretch
  \compacttablecols
  \setlength{\tabcolsep}{6pt}
  \caption{\textbf{Model-averaged satellite for the systematic factor $Z$.}}
  \label{tab:regZ}
  \resizebox{0.72\textwidth}{!}{%
    \begin{tabular}{lccc}
      \toprule
      & Post. mean & Post. SD & Incl. prob. \\
      \midrule
      Intercept
      & 12.827 & 0.636 & 1.00 \\
      $\log(\text{Investment pc})_{t}$
      & $0.039^{***}$ & 0.010 & 1.00 \\
      $\log(\text{Investment pc})_{t-2}$
      & 0.016 & 0.014 & 0.64 \\
      $\log(\text{Oil real})_{t-4}$
      & $-0.004^{***}$ & 0.003 & 0.63 \\
      Inflation YoY$_{t-4}$
      & 0.059 & 0.059 & 0.52 \\
      $\log(\text{Oil real})_{t}$
      & $-0.001^{***}$ & 0.002 & 0.39 \\
      $\log(\text{Oil real})_{t-1}$
      & $-0.002^{***}$ & 0.003 & 0.33 \\
      $\log(\text{GDP pc})_{t-4}$
      & $-0.015^{***}$ & 0.023 & 0.32 \\
      $\log(\text{GDP pc})_{t}$
      & $-0.017^{***}$ & 0.026 & 0.30 \\
      \midrule
      Observations
      & \multicolumn{3}{r}{152} \\
      $R^2$ / adj. $R^2$
      & \multicolumn{3}{r}{0.89 / 0.88} \\
      Out-of-sample $R^2$
      & \multicolumn{3}{r}{0.85} \\
      \bottomrule
    \end{tabular}%
  }
  \par
  \medskip
  \begin{minipage}{0.72\textwidth}
  \footnotesize
  \noindent \textbf{Note:} The table reports the posterior mean, posterior
  standard deviation, and inclusion probability of each satellite coefficient
  under Schwarz model averaging. Stars denote a posterior probability of the
  coefficient sign exceeding $0.95$, $0.975$, and $0.99$, respectively.
  Because the reported quantities are moments of the model-averaging mixture,
  stars on terms whose inclusion probability is well below one reflect the sign
  probability \emph{conditional on inclusion} and should not be read as
  unconditional evidence of relevance. $R^2$ and out-of-sample $R^2$ are
  computed for the single best specification.
  \end{minipage}
\end{table}
\subsection{Dynamic responses}
\label{sec:empirical_results}
We propagate the one-standard-deviation innovation in the GPR equation through
the estimated system: first to the macro-financial variables, then to the
systematic factor, and finally to portfolio PDs. All responses are reported as
posterior medians with pointwise $68\%$ and $90\%$ credible bands.
\textit{Macro-financial responses.}
Figure~\ref{fig:irfs_var} reports the responses of the macro-financial
variables. The GPR innovation depresses real activity: investment, GDP per
capita, and private employment per capita fall and stay below baseline for
several quarters, while the real oil price and inflation respond more
transiently. The pattern is consistent with the contractionary effects of
geopolitical risk documented by \citet{caldara2022geopolitical}.
\begin{figure}[!htbp]
    \centering
    \includegraphics[width=\linewidth]{figures/main/graph_var_irf.pdf}
    \caption{\textbf{Macro-financial responses to a GPR innovation.}
    \textit{Notes: Posterior medians (solid) with $68\%$ (dark) and $90\%$
    (light) pointwise credible bands. Responses are generalized impulse
    responses to a one-standard-deviation innovation in the GPR equation.}}
    \label{fig:irfs_var}
\end{figure}
\medskip
\textit{Transmission to the systematic factor and portfolio PDs.}
Figure~\ref{fig:girf_Z_PD} reports the transmission of the innovation to the
systematic factor and to portfolio PDs. Panel~(a) shows that the innovation
lowers the systematic factor: the response is negligible on impact, turns
negative within a few quarters, and reverts toward baseline thereafter, with
bands widening at longer horizons. This path combines the macro-financial
propagation in the VAR with the positive exposure of $Z$ to real activity in
the satellite equation. Through the Merton--Vasicek mapping, panel~(b)
translates this movement into an increase in portfolio PDs. The median PD
response peaks after three quarters at $+0.033$ percentage points, about
$1\%$ of the through-the-cycle level; the posterior probability of a positive
response at the peak is $0.94$. The median response cumulates to about
$0.26$ percentage-point-quarters of additional default risk over the
twelve-quarter horizon.
\begin{figure}[!htbp]
  \centering
  \begin{subfigure}[t]{0.49\linewidth}
    \centering
    \includegraphics[width=\linewidth]{figures/main/graph_z.pdf}
    \caption{Systematic factor $Z$ (s.d.).}
    \label{fig:girf_Z}
  \end{subfigure}
  \hfill
  \begin{subfigure}[t]{0.49\linewidth}
    \centering
    \includegraphics[width=\linewidth]{figures/main/graph_pd.pdf}
    \caption{Portfolio PD (percentage points).}
    \label{fig:girf_pd}
  \end{subfigure}
  \caption{\textbf{Generalized responses to a GPR innovation: $Z$ and PD.}
  \textit{Notes: Panel (a) reports the response of the systematic credit
  factor $Z$ implied by the VAR and the model-averaged satellite. Panel (b)
  maps this response into portfolio default probabilities through the
  Merton--Vasicek mapping with $\widehat p=3.20\%$ and $\widehat\rho=0.051$.
  Posterior medians (solid) with $68\%$ (dark) and $90\%$ (light) pointwise
  credible bands. Responses are to a one-standard-deviation innovation in the
  GPR equation.}}
  \label{fig:girf_Z_PD}
\end{figure}
\textit{Tail responses.}
Capital and provisioning depend on the upper tail of the conditional
default-probability distribution rather than on its center.
Corollary~\ref{cor:pd_ar} delivers this tail response in closed form.
Figure~\ref{fig:girf_pd_ar} reports the generalized response of the portfolio
PD-at-Risk (PD-aR). Its timing matches the mean response, with a peak at the
third quarter, but its magnitude is larger: the $99\%$ PD-aR rises by
$+0.050$ percentage points at the peak and the $99.9\%$ PD-aR by $+0.056$,
against $+0.033$ for the mean, a ratio of about $1.5$. This amplification
follows from the curvature of the Merton--Vasicek mapping, under which a given
decline in the systematic factor shifts the upper tail of the conditional PD
distribution more than its center. Measured against the through-the-cycle
PD-aR levels of $8.6\%$ ($99\%$) and $11.8\%$ ($99.9\%$), the response is
modest, but it indicates that a geopolitical innovation deteriorates credit
quality more in the tail than at the center of the distribution.
\begin{figure}[!htbp]
  \centering
  \begin{subfigure}[t]{0.49\linewidth}
    \centering
    \includegraphics[width=\linewidth]{figures/main/pd_var0p99_dralacbn_real_side_1sd.pdf}
    \caption{Response of the $99\%$ PD-aR.}
    \label{fig:girf_pd_ar_a}
  \end{subfigure}
  \hfill
  \begin{subfigure}[t]{0.49\linewidth}
    \centering
    \includegraphics[width=\linewidth]{figures/main/pd_mean_vs_var_dralacbn_real_side_1sd.pdf}
    \caption{Mean response versus PD-aR responses.}
    \label{fig:girf_pd_ar_b}
  \end{subfigure}
  \caption{\textbf{Tail responses of portfolio default probabilities to a GPR
  innovation.}
  \textit{Notes: Panel (a) reports the generalized response of the $99\%$
  portfolio PD-at-Risk defined in Corollary~\ref{cor:pd_ar}, with
  posterior median (solid) and $68\%$ (dark) and $90\%$ (light) pointwise
  credible bands. Panel (b) compares the posterior-median responses of the
  mean PD and of the $99\%$ and $99.9\%$ PD-aR. All responses are to a
  one-standard-deviation innovation in the GPR equation, with
  $\widehat p=3.20\%$ and $\widehat\rho=0.051$.}}
  \label{fig:girf_pd_ar}
\end{figure}
\subsection{Geopolitical episodes and state dependence}
\label{sec:scenario_2001}
Stress-testing frameworks require severe yet plausible scenarios. The
estimated VAR provides a direct device for constructing them: its history of
standardized reduced-form GPR innovations identifies large historical shocks
in the units of the model (Figure~\ref{fig:gpr_innov}). The three largest
positive innovations in the sample correspond to recognizable episodes: the
invasion of Kuwait in 1990:Q3 ($3.14$ standard deviations), the September 11
attacks in 2001:Q3 ($4.45$), and the invasion of Ukraine in 2022:Q1
($3.54$). Because the response of Proposition~\ref{prop:GIRF-PD} is
history-dependent through the conditional mean of the systematic factor, each
episode can be evaluated at its own initial conditions, and these episodes
happen to span the credit cycle: the Gulf War struck a stressed credit system
(model-implied PD of $4.6\%$ at the shock date), the September 11 attacks an
intermediate one ($2.0\%$), and the Ukraine invasion a benign one ($1.4\%$).
The macro-financial propagation is identical across these experiments, since
the VAR is linear; the comparison isolates the position of the credit cycle
in the Merton--Vasicek map. Parameters are full-sample estimates, so the
exercise is a counterfactual within the estimated model rather than a
real-time analysis.
\begin{figure}[!htbp]
    \centering
    \includegraphics[width=0.80\linewidth]{figures/main/graph_e1.pdf}
    \caption{\textbf{Standardized GPR innovations.}
    \textit{Notes: Standardized reduced-form innovations in the GPR equation,
    $\varepsilon_{\mathrm{GPR},t}=u_{1,t}/\sqrt{\Sigma_{1,1}}$, where
    $u_{1,t}$ is the innovation to the GPR equation of the Bayesian VAR and
    $\Sigma_{1,1}$ is its innovation variance, evaluated at posterior-median
    parameters. Values are in standard-deviation units; positive values
    correspond to unexpected increases in geopolitical risk.}}
    \label{fig:gpr_innov}
\end{figure}
Table~\ref{tab:state_dependence} reports the responses. The historical
comparison delivers the central result, displayed in panel~(a) of
Figure~\ref{fig:state_dependence_main}: evaluated at its own date, the Gulf
War innovation produces a larger PD response ($+0.278$ percentage points at
the peak) than the September 11 innovation ($+0.219$), although the latter
shock is more than forty percent larger. A smaller geopolitical shock
generated a larger credit-risk response because it struck a weaker credit
system. The marginal credible bands of the two responses overlap, but the
comparison is sharper than this overlap suggests: the two experiments are
evaluated on identical posterior draws, so that estimation uncertainty is
common to both, and the posterior probability that the Gulf War response
exceeds the September 11 response at the peak is $0.94$. Panel~(b)
generalizes the comparison beyond the three episodes: for every admissible
history in the sample, it plots the peak response to a common
one-standard-deviation innovation against the baseline expected PD at the
shock date, at posterior-median parameters. The peak response rises from
$0.029$ percentage points at the most benign state (2023:Q4) to $0.131$ at
the most stressed (2009:Q4), a factor of $4.6$ over the historically observed
range of credit conditions.
\begin{figure}[!htbp]
  \centering
  \begin{subfigure}[t]{0.49\linewidth}
    \centering
    \includegraphics[width=\linewidth]{figures/main/fig_state_dependence_inversion.pdf}
    \caption{The two largest episodes, own shock sizes.}
    \label{fig:state_dependence_inversion}
  \end{subfigure}
  \hfill
  \begin{subfigure}[t]{0.49\linewidth}
    \centering
    \includegraphics[width=\linewidth]{figures/main/fig_state_dependence_curve.pdf}
    \caption{Peak response across all histories.}
    \label{fig:state_dependence_curve}
  \end{subfigure}
  \caption{\textbf{Geopolitical episodes and state dependence of the PD
  response.}
  \textit{Notes: Panel (a) replays each episode at its own date with its own
  historical shock size (Gulf War: $3.14$ s.d.; September 11: $4.45$ s.d.);
  posterior medians with $90\%$ pointwise credible bands on identical
  posterior draws. Panel (b) plots the peak response to a common
  one-standard-deviation GPR innovation against the baseline expected PD at
  the shock date, for every admissible history in the sample, at
  posterior-median parameters; colored points mark the three episodes and the
  end-of-sample state. A linear specification of the default rate would imply
  a horizontal line in panel (b).}}
  \label{fig:state_dependence_main}
\end{figure}
These results characterize the state dependence of the macro-to-credit
transmission. The state of the credit cycle is the operative margin: the same
one-standard-deviation GPR innovation generates a larger PD response when
initial credit conditions are weaker. By contrast, at a given initial state,
the size of the shock scales the response almost proportionally. From the
end-of-sample state, for instance, a September-11-sized shock yields
$+0.155$ percentage points, which is $4.6$ times the one-standard-deviation
response for a $4.45$-standard-deviation shock. The timing of the peak
response, at three quarters, and the posterior probability of a positive
response, equal to $0.94$, are unchanged across states. State dependence
therefore operates through the Merton--Vasicek credit mapping, rather than
through the macro-financial dynamics.
The same evidence rationalizes the two-block architecture. A linear VAR that
included the default rate directly would imply state-independent responses, a
horizontal line in panel~(b) of Figure~\ref{fig:state_dependence_main}, and
proportionality between the Gulf War and September 11 responses in
panel~(a). Instead, the estimated response is strongly state-dependent: even
though the relation between baseline expected PD and the peak response is close
to linear over the observed range, its positive slope is economically large,
and the Gulf War/September 11 comparison displays an inversion.
\begin{table}[!htbp]
  \centering
  \scriptsize
  \tablestretch
  \compacttablecols
  \setlength{\tabcolsep}{4pt}
  \caption{\textbf{Geopolitical episodes evaluated at their own initial conditions.}}
  \label{tab:state_dependence}
  \resizebox{1\textwidth}{!}{%
    \begin{tabular}{lccccc}
      \toprule
      Episode
      & \makecell{PD at shock\\date (pp)}
      & \makecell{Shock\\(s.d.)}
      & \makecell{Peak $\Delta$PD,\\own shock (pp)}
      & \makecell{Peak $\Delta$PD,\\1 s.d. (pp)}
      & Ratio \\
      \midrule
      Gulf War (1990:Q3)
      & 4.59 & 3.14 & \makecell{$0.278$\\{\scriptsize $[0.099,\,0.477]$}} & 0.087 & 2.59 \\
      September 11 (2001:Q3)
      & 1.97 & 4.45 & \makecell{$0.219$\\{\scriptsize $[0.077,\,0.376]$}} & 0.047 & 1.41 \\
      Ukraine invasion (2022:Q1)
      & 1.36 & 3.54 & \makecell{$0.106$\\{\scriptsize $[0.059,\,0.162]$}} & 0.029 & 0.87 \\
      End of sample (2025:Q1)
      & 1.45 & -- & -- & 0.033 & 1.00 \\
      \bottomrule
    \end{tabular}%
  }
  \par
  \medskip
  \begin{minipage}{\textwidth}
  \footnotesize
  \noindent \textbf{Note:} Generalized responses of the portfolio PD are
  evaluated conditional on the history preceding each episode, on identical
  posterior draws. ``PD at shock date'' is the posterior median of the
  model-implied baseline PD at $h=0$. ``Own shock'' replays the episode's
  standardized reduced-form GPR innovation; brackets report $68\%$ credible
  intervals. ``1 s.d.'' applies a common one-standard-deviation innovation to
  isolate state dependence. ``Ratio'' is the corresponding peak relative to the
  end-of-sample state. All peaks occur at $h=3$. The episodes are the three
  largest positive GPR innovations in the sample. Parameters are full-sample
  estimates.
  \end{minipage}
\end{table}
\subsection{The cost of perfect foresight}
\label{sec:perfect_foresight}
The closed-form expression allows us to quantify the effect of predictive
uncertainty separately from estimation uncertainty. This distinction matters
because supervisory stress tests typically treat the macroeconomic scenario
as a known path: banks ``should assume the subsequent path of a variable to
be known in line with the scenario,'' and provisions are projected under one
path per scenario \citep[Box~1, paragraphs~130 and~136]{EBA2025method}. In
our framework, this convention corresponds to suppressing part of the
conditional variance that enters the Merton--Vasicek map. We therefore
recompute the PD projections on the same posterior draws under three
variants: the exact closed form of Proposition~\ref{prop:GIRF-PD}; a
perfect-foresight variant, which treats the macro-financial path as known but
retains the satellite-error variance, the closest counterpart to the
supervisory convention; and a plug-in variant, which evaluates the
Merton--Vasicek map at the conditional-mean path. Since the posterior draws
are identical across variants, differences are due only to predictive
uncertainty.
Table~\ref{tab:pf_bias} and Figure~\ref{fig:pf} show that perfect foresight
mainly affects projected PD levels, not impulse responses. Because the
Merton--Vasicek map is convex in the relevant region, suppressing predictive
uncertainty understates expected PDs. The perfect-foresight variant
understates the exact PD path by $2.1\%$ at the one-year horizon and by
$6.7\%$ at the three-year horizon. The plug-in approximation produces larger
biases, equal to $3.9\%$ and $8.4\%$, respectively. For illustration, on a
EUR~10 billion corporate credit portfolio with a $45\%$ loss-given-default,
the three-year plug-in bias corresponds to approximately EUR~6 million of
expected-loss understatement. By contrast, impulse responses are almost
unchanged: the convexity correction enters both the shocked and baseline
evaluations and largely cancels when the response is computed as a
difference, leaving peak-response biases at most about $2\%$ for both the
one-standard-deviation and the September-11-sized shock. Perfect foresight
therefore has limited effects on the estimated sensitivity of credit risk to
shocks, but it lowers projected default-risk levels, which are directly
relevant for provisioning and capital projections. By convexity, the
understatement increases with the horizon and in adverse states such as those
documented in Section~\ref{sec:scenario_2001}.
\begin{table}[!htbp]
  \centering
  \scriptsize
  \tablestretch
  \compacttablecols
  \setlength{\tabcolsep}{6pt}
  \caption{\textbf{The cost of perfect foresight: PD levels and responses.}}
  \label{tab:pf_bias}
  \resizebox{0.96\textwidth}{!}{%
    \begin{tabular}{llccccc}
      \toprule
      & & Exact & PF & Plug-in & Bias PF (\%) & Bias plug-in (\%) \\
      \midrule
      PD level (pp) & $h=4$
      & 1.427 & 1.397 & 1.371 & $-2.1$ & $-3.9$ \\
      PD level (pp) & $h=8$
      & 1.458 & 1.390 & 1.366 & $-4.7$ & $-6.3$ \\
      PD level (pp) & $h=12$
      & 1.503 & 1.401 & 1.377 & $-6.7$ & $-8.4$ \\
      \midrule
      Peak $\Delta$PD (pp), one s.d. & $h=3$
      & 0.033 & 0.033 & 0.033 & $-0.1$ & $-1.2$ \\
      Peak $\Delta$PD (pp), 2001:Q3 & $h=3$
      & 0.155 & 0.154 & 0.152 & $-0.8$ & $-2.0$ \\
      \bottomrule
    \end{tabular}%
  }
  \par
  \medskip
  \begin{minipage}{0.96\textwidth}
  \footnotesize
  \noindent \textbf{Note:} Posterior medians are computed on identical posterior
  draws. The variants differ only in the predictive variance entering the closed
  forms of Section~\ref{sec:framework}. ``Exact'' is
  Proposition~\ref{prop:GIRF-PD}. ``PF'' treats the macro-financial path as known
  but retains the satellite-error variance, which is the closest counterpart to
  the perfect-foresight convention of supervisory stress tests
  \citep{EBA2025method}. ``Plug-in'' evaluates the Merton--Vasicek map at the
  conditional-mean path. Biases are relative to the exact variant. Negative level
  biases indicate understatement of the expected PD path.
  \end{minipage}
\end{table}
\begin{figure}[H]
  \centering
  \includegraphics[width=0.72\linewidth]{figures/main/fig_perfect_foresight_shocked_levels.pdf}
  \caption{\textbf{The cost of perfect foresight.}
  \textit{Notes: The figure compares the PD path along the September~11
  (2001:Q3) scenario under the exact closed form (with a $68\%$ posterior band)
  with the paths obtained when the macro-financial trajectory is treated as
  known (PF) and when the Merton--Vasicek map is evaluated at the
  conditional-mean path (plug-in). All variants are computed on identical
  posterior draws; the corresponding biases of the impulse responses are
  reported in Table~\ref{tab:pf_bias}.}}
  \label{fig:pf}
\end{figure}
\subsection{Using shorter default histories}
\label{sec:short_default_histories}
A key motivation for the modular architecture is that macro-financial variables
are often observed over long samples, whereas default histories are shorter and
may be affected by changes in definitions, regulation, accounting standards, or
portfolio composition. In such settings, a fully joint macro-credit model would
either discard part of the macro-financial history or require a homogeneous
default series over the entire macro sample. By contrast, the proposed
framework estimates the VAR on the full macro-financial sample and re-estimates
only the credit-risk bridge on the available default window.
To illustrate this feature, we keep the macro-financial VAR estimated over
1986:Q1--2024:Q4 and re-estimate the Merton--Vasicek inversion and the
satellite equation on shorter default samples. Figure~\ref{fig:short_default_windows}
reports the resulting portfolio-PD responses when the credit-risk block starts
in 2005:Q1 and in 2015:Q1. The same one-standard-deviation GPR innovation is
used in all cases, so differences across panels reflect the information
contained in the shorter default histories and the corresponding satellite
bridge, not changes in the macro-financial propagation mechanism.
The responses remain positive in both exercises, but their magnitude and
uncertainty change with the length of the default sample. The 2005 window still
contains the global financial crisis and therefore includes a large stressed
credit episode from which the satellite can learn the sensitivity of default
risk to macro-financial conditions. The 2015 window is much shorter and mostly
covers a benign credit environment, which weakens the empirical identification
of the credit-risk bridge and widens posterior uncertainty. This comparison
highlights the role of the modular design: the long macro-financial sample can
still be used to identify and propagate the geopolitical-risk shock, while the
credit-risk block can be adapted to the default history that is actually
available for the portfolio under study.
\begin{figure}[!htbp]
    \centering
    \begin{subfigure}[t]{0.48\textwidth}
        \centering
        \includegraphics[width=\textwidth]{figures/main/pd_girf_window_2005_1sd.pdf}
        \caption{Default sample starts in 2005:Q1.}
        \label{fig:short_default_window_2005}
    \end{subfigure}
    \hfill
    \begin{subfigure}[t]{0.48\textwidth}
        \centering
        \includegraphics[width=\textwidth]{figures/main/pd_girf_window_2015_1sd.pdf}
        \caption{Default sample starts in 2015:Q1.}
        \label{fig:short_default_window_2015}
    \end{subfigure}
    \caption{\textbf{Portfolio-PD responses with shorter default histories.}
    \textit{Notes: The macro-financial VAR is estimated on the full
    1986:Q1--2024:Q4 sample. The Merton--Vasicek credit-risk block and the
    satellite equation are re-estimated on shorter default samples starting in
    2005:Q1 and 2015:Q1, respectively. Each panel reports the generalized
    response of portfolio default probabilities to a one-standard-deviation
    innovation in the GPR equation, with posterior medians and pointwise
    credible bands.}}
    \label{fig:short_default_windows}
\end{figure}
